\section{The remaining integral dimension-forty-eight exclusions}
\label{sec:d48-exclusions}
\begin{theorem}
The based rings $R9\_1\_62$ and $R9\_1\_87$ admit no fusion-category
categorification over $\mathbb C$. The same argument distinguishes them from
$R9\_1\_12$ and $R9\_1\_108$, for which unitary constructions are supplied
separately.
\end{theorem}

\paragraph{First reduction.}
All four rings have simple dimensions $(1^4,2^3,4^2)$, invertible group
$V_4$, and trivial universal grading group. The latter assertion is checked
by generating the adjoint subring from all products $XX^*$.
A categorification is solvable by the Etingof--Nikshych--Ostrik
$p^a q^b$ theorem \cite[Theorem~1.6]{ENOSolvable}. In a sequence of cyclic
prime-order extensions and equivariantizations constructing it
(\cite[Proposition~4.4]{ENOSolvable}), the final step cannot be a nontrivial
extension, since its universal grading is trivial. It must therefore be a
$C_p$-equivariantization. Its canonical $\operatorname{Rep}(C_p)$ embeds
centrally, so $p$ divides the order-four invertible group and $p=2$.
Write $\mathcal C=\mathcal D^{C_2}$, and let $e$ be the corresponding sign
object of $\mathcal C$.

The forgetful functor is completely constrained on simple-object orbits.
A pair $\{X,eX\}$ forgets to one simple of dimension $\operatorname{FPdim}(X)$;
a simple fixed by tensoring with $e$ forgets to two distinct simples,
exchanged by the action, of dimension $\operatorname{FPdim}(X)/2$.
The projective stabilizers are cyclic and have trivial complex Schur
multiplier, so no further projective-representation dimension occurs.
The exact three possibilities are as follows; the entries give the
zero-based input label of $e$.
\[
\begin{array}{c|ccc}
\text{input}&(1^8,4)&(1^4,2,4)&(1^4,2^5)\\\hline
R9\_1\_12&7&1&2\\
R9\_1\_62&5&6&7\\
R9\_1\_87&5&4&3\\
R9\_1\_108&5&7&6
\end{array}
\]
The quotient has dimension~24 in each case. The integer verifier
\nolinkurl{verify_d48_exclusions.py} checks these orbits and every dimension
identity directly against the input tensors.

\paragraph{The near-group branch.}
For profile $(1^8,4)$, writing $\rho$ for the unique noninvertible gives
$g\rho=\rho g=\rho$ and
$\rho^2=\sum_{g\in G}g+2\rho$. The acting involution fixes exactly two
invertibles. The complete near-group action theorem of
Section~\ref{sec:ng24-actions} applies: its equivariantization is
$R9\_1\_12$ for $G=D_8$, or $R9\_1\_108$ for $G=Q_8$.
In particular this branch excludes the two rings in the theorem.
This is an exhaustive categorical statement, including all six associator
classes and all tensorators, not merely a test of strict group actions.

\paragraph[The rank-nine quotient branch]{The $(1^4,2^5)$ branch.}
The involution on the quotient fixes exactly two invertibles and one
dimension-two simple $m$. It exchanges the remaining invertibles and pairs
up the remaining four dimension-two simples. The fixed $m$ is self-dual.
The source products of either of its two equivariant lifts force
\[
 m^2=\sum_{g\in\operatorname{Inv}(\mathcal D)}g.
\]
The exact source products used here, in the input labeling, are
$3\cdot3=0+7+8$ for input~12, $1\cdot2=0+5+8$ for input~62,
$6\cdot6=0+5+7$ for input~87, and $1\cdot2=0+5+8$ for input~108.
The last summand in each product forgets to the two nonfixed invertibles.

The quotient ring is necessarily multiplicity-free. For a nonfixed
quotient simple $Z$, its coefficient in the forgetful image of a unique
source simple is one. Thus the tensor-functor identity and the
multiplicity-free source tensor bound every $N_{XY}^Z$ by one.
Invertible output coefficients are bounded by one by rigidity.
The only remaining possible output is $m$. If an input is invertible the
bound is immediate. Otherwise, if $Y\ne m$, Frobenius reciprocity gives
$N_{XY}^{m}=N_{X^*m}^{Y}\le1$, since $Y$ is nonfixed.
For $Y=m$, reciprocity reduces the coefficient to one in $m^2$, which is
already known. This proves the claim without assuming that a quotient of
a multiplicity-free category is generally multiplicity-free.

A direct binary-tensor enumeration now completes the quotient ring,
without consulting a catalog of quotient rings. Label the quotient so that
its dimensions are $(1,1,1,1,2,2,2,2,2)$ and its involution is
$(2\ 3)(5\ 6)(7\ 8)$. The source sign orbits determine the based
restriction matrix $F$, up to this harmless choice of labels.
The verifier \nolinkurl{direct_quotient.py} enumerates every duality involution
compatible with $F$, dimensions, and the action. It identifies tensor
coefficients only under Frobenius reciprocity and the action, then
exhausts all binary assignments satisfying unit, duality, dimension, and
all identities
\[
 \sum_h N_{ij}^{h}(\mathcal C) F_{hk}
 =\sum_{a,b}F_{ia}F_{jb}N_{ab}^{k}(\mathcal D).
\]
For every variable, the positive linear equations also supply an explicit
upper-bound witness at most one, independently confirming that no larger
nonnegative integer coefficient is omitted. Every surviving assignment
is checked against every associativity identity.
There are respectively $8,8,4,4$ possible dualities for sources
$12,62,87,108$; the numbers of linearly feasible completions are
$36,36,4,4$. Exactly four labeled associative completions remain for
each source. Only after this exhaustive completion does the verifier
compare them with $R9\_1\_117$, and all sixteen comparisons give explicit
based isomorphisms. Thus $R9\_1\_117$ is the only possible quotient,
independently of completeness of any input census.

The complete triality-action classification, established in
\nolinkurl{work/resolve_integral/TRIALITY_ACTIONS.md}, now applies.
Every categorification of $R9\_1\_117$ has neutral category
$\operatorname{Rep}(D_8)$ in its canonical $C_3$ grading, and every compatible
$C_2$ action gives only $R9\_1\_12$ or $R9\_1\_108$.
The argument includes the neutral-category restriction, all extension
associator parameters, and the two possible coherence classes.
Consequently this branch also excludes~62 and~87.

\paragraph{The rank-six branch.}
Here write $G$ for the order-four invertible group and $m,n$ for the
unique simples of dimensions~2 and~4. No assumption that $G$ is $V_4$ is
needed at this stage. Dimension, rigidity, and Frobenius reciprocity force
\[
 gm=mg=m,\quad gn=ng=n,\quad
 m^2=\sum_G g,\quad mn=nm=2n,\quad
 n^2=\sum_G g+2m+2n.
\]
In particular this category, denoted $\mathcal B$, has trivial universal
grading. Solvability once again makes it a $C_2$-equivariantization
$\mathcal B=\mathcal E^{C_2}$. The orbit calculation gives
$\operatorname{FPdim}(\mathcal E)=12$, profile $(1^4,2^2)$; the involution
exchanges its two dimension-two simples and transposes two of its
invertibles.

We give the necessary ring reduction explicitly. Let $A$ be the
order-four group of invertibles of $\mathcal E$, and let $X,Y$ be its
dimension-two simples. Duality shows that the left and right actions of
$A$ on $\{X,Y\}$ coincide. If this action were transitive, its stabilizer
$H$ would have order~2 and
$XX^*=\sum_Hh+Z$ for one dimension-two simple $Z$. This expression is
self-dual, so $Z$ is self-dual and hence both $X,Y$ are self-dual.
Choose $g\notin H$; then $Y=gX=Xg$ and $g^2\in H$, whence
$Y^2=g^2X^2=X^2$. An automorphism exchanging $X,Y$ would have to preserve
the unique noninvertible summand $Z$ of this common square, a contradiction.
Thus $A$ fixes both simples, and $XX^*=YY^*=\sum_Aa$.
If both were self-dual, write $XY=aX+bY$; the dimension identity gives
$a+b=2$, whereas $X(XY)=X^2Y=4Y$ forces $a=0$, and
$(XY)Y=XY^2=4X$ then gives an incompatible invertible summand.
Therefore $X^*=Y$, and associativity, or Frobenius reciprocity, gives
\[
 XY=YX=\sum_Aa,\qquad X^2=2Y,\qquad Y^2=2X.
\]
This is the canonical faithful $C_3$ grading with neutral pointed category
on $A$ and one simple in each nonneutral component.

A rank-one homogeneous module forces the neutral associator to be
cohomologically trivial, so the neutral category is $\operatorname{Vec}_A$.
The Brauer--Picard action of a generator of $C_3$ is an isometry
$M\in O(A\oplus A^*,q)$, $q(a,\chi)=\chi(a)$, whose block
$\beta:A^*\to A$ is invertible. These standard extension facts are also
stated and proved in Jordan--Larson,
\cite[Lemma~2.6 and Proposition~3.1]{JordanLarson}. If $A=C_4$, its split quadratic group has exactly four
automorphisms: $I,-I$, coordinate interchange, and its negative. None has
order~3, so this case is impossible. The verifier exhausts all~256 candidate
matrices modulo~4 to check this assertion.

For $A=V_4$, identify $A\oplus A^*$ with $\mathbb F_2^4$, with
$q(a_1,a_2,\chi_1,\chi_2)=a_1\chi_1+a_2\chi_2$ in exponent notation.
The split orthogonal group has~72 elements. Exactly six of them have order~3
and invertible $\beta$. In every one of these six cases,
\begin{equation}\label{eq:d12-no-isotropic-fixed}
 Mv=v,\quad q(v)=0\quad\Longrightarrow\quad v=0.
\end{equation}
This is checked by exhaustive exact enumeration in the same verifier.
Equivalently, the fixed space is zero or is an anisotropic plane.
No choice of scalar associator or of a lift of the action changes this
necessary fixed-vector statement. Jordan--Larson's six-category
classification \cite[Theorem~1.4]{JordanLarson} provides an independent description of the same two matrix
orbits and their three $C_3$ associator twists.

\paragraph{A characteristic central sign.}
Put $\mathcal K=\mathcal E\rtimes C_2$. It is Morita equivalent to
$\mathcal B$ and is faithfully $S_3$-graded with neutral category
$\operatorname{Vec}_{V_4}$, because the involution exchanges the two
nonneutral $C_3$ components. Any central invertible of $\mathcal K$ has
degree in $Z(S_3)=1$; restriction of its half-braiding to the neutral
category gives an $M$-fixed vector in $A\oplus A^*$.
Its twist is precisely $(-1)^{q(v)}$ for that vector $v$. By
\eqref{eq:d12-no-isotropic-fixed}, every central invertible of twist $+1$
has trivial restriction. The kernel of this restriction consists exactly
of grading characters $\operatorname{Hom}(S_3,\mathbb C^*)=C_2$.
Indeed its underlying object is the unit, its half-braiding is trivial
on the neutral component, and its scalar is constant on each indecomposable
homogeneous bimodule component and multiplicative in the grading.
Both kernel objects have twist $+1$. Hence $\mathcal K$, and therefore
$\mathcal B$, has exactly one nontrivial central invertible of twist $+1$.
This proof does not require that any nonzero invariant fermion lift to the
center.

All categories under consideration are weakly integral and therefore have
the canonical pseudo-unitary spherical structure. The braided equivalence
of centers induced by Morita equivalence preserves its twist.
Moreover $U(\mathcal B)=1$, so the forgetful map
$\operatorname{Inv}Z(\mathcal B)\to\operatorname{Inv}(\mathcal B)$ is injective.
The unique nontrivial boson thus has a nontrivial underlying object and
generates a central Tannakian $\operatorname{Rep}(C_2)$.
Every tensor autoequivalence preserves this characteristic central
subcategory.

Return to the original equivariantization
$\mathcal C=\mathcal B^{C_2}$. Equivariantizing this stable central
$\operatorname{Rep}(C_2)$ gives a central Tannakian subcategory of dimension~4
in $\mathcal C$: the equivariantization remains symmetric, and the
composite forgetful functor to $\operatorname{Rep}(C_2)$ and then to
vector spaces is a symmetric fiber functor. It is $\operatorname{Rep}(V_4)$ or
$\operatorname{Rep}(C_4)$. The latter is impossible since
$\operatorname{Inv}(\mathcal C)=V_4$. Thus all four invertible objects of
$\mathcal C$ belong to central $\operatorname{Rep}(V_4)$.
We may therefore choose any of its three signs for de-equivariantization,
in particular the sign in the $(1^8,4)$ column of the displayed table.
The already exhaustive near-group branch forces $R9\_1\_12$ or
$R9\_1\_108$, excluding~62 and~87 in the last remaining branch. This
completes the proof.

\paragraph{Replay and scope.}
The exact finite receipt is
\nolinkurl{work/resolve_rank9_integral/d48_exclusions_receipt.json}.
Run \nolinkurl{python work/resolve_rank9_integral/verify_d48_exclusions.py}
for the dimension identities, all source orbit profiles, the independent
replay of all completed quotient tensors, the $C_4$ obstruction, and all six
quadratic fixed-space checks. The detailed near-group theorem and its
independent matrix/tensor replay are
\nolinkurl{work/resolve_rank9_neargroup/NG_ACTIONS.tex} and
\nolinkurl{verify_ng_actions.py}. The triality action classification is
\nolinkurl{work/resolve_integral/TRIALITY_ACTIONS.md}.
The exhaustive quotient completion and its receipt are
\nolinkurl{work/resolve_rank9_neargroup/direct_quotient.py} and
\nolinkurl{direct_quotient.json}. These categorical reductions and exact
finite enumerations supply the complete exclusion;
no infeasible numerical search or failure to find a model is used.
